% Figaro Paper B1: From Firms to Transaction-Scoped Institutions
% A Coasean Re-Examination Under Costless Bilateral Enforcement
% (Institutional-economics half of the former Paper B; the political-economy
% argument is now in figaro3b2.tex.)
% Preprint, April 2026

\documentclass[11pt,a4paper]{article}

\usepackage[utf8]{inputenc}
\usepackage[T1]{fontenc}
\usepackage{amsmath,amssymb,amsthm}
\usepackage{mathtools}
\usepackage{booktabs}
\usepackage{graphicx}
\usepackage{hyperref}
\usepackage{cleveref}
\usepackage{enumitem}
\usepackage{xcolor}
\usepackage[margin=1in]{geometry}
\usepackage{natbib}
\bibliographystyle{plainnat}

\newtheorem{theorem}{Theorem}[section]
\newtheorem{proposition}[theorem]{Proposition}
\newtheorem{corollary}[theorem]{Corollary}
\theoremstyle{definition}
\newtheorem{definition}[theorem]{Definition}
\theoremstyle{remark}
\newtheorem{remark}[theorem]{Remark}

\newcommand{\figaro}{\textsc{Figaro}}
\newcommand{\fig}{\textsc{Fig}}
\newcommand{\pay}{P}
\newcommand{\cumval}{G}

\title{%
  From Firms to Transaction-Scoped Institutions:\\
  A Coasean Re-Examination Under Costless Bilateral Enforcement%
}

\author{
  Alessandro Daliana\thanks{Corresponding author. \texttt{adaliana@figaro.org}}
}

\date{April 2026}

\begin{document}

\maketitle

% ════════════════════════════════════════════════════════════════════
\begin{abstract}
\citet{coase1937} argued that firms exist because the cost of
discovering prices, negotiating terms, and enforcing contracts in
open markets sometimes exceeds the cost of hierarchical management.
\citet{williamson1985} refined this into Transaction Cost Economics
(TCE), identifying asset specificity, uncertainty, and frequency as
drivers of vertical integration. \citet{hart1995} located the firm
boundary at the allocation of residual control rights over
non-contractible decisions. Across these traditions, a common premise
is that contract enforcement between strangers is costly, slow, and
jurisdiction-dependent. This paper examines what happens to the
boundary of the firm when one of those underlying costs collapses to
a fixed, jurisdiction-displaced capital lockup.

We take as given the existence of a settlement
primitive---developed formally in its mechanism-design treatment
and implemented as a verified Solidity contract---that prices
the cost of trust at exactly $2\times$ the transaction value in
temporarily locked
capital, with no recurring extraction, no arbitrator, and no
jurisdictional dependency. Treating this primitive as exogenous, we
ask two questions. First, what does the existence of such a
primitive imply for the Coasean threshold that defines the boundary
of the firm? Second, what is the firm's status in the organizational
space once that threshold has shifted?

Our argument is that the firm has always been a coordination
mechanism with a particular layer added --- \emph{subordination}:
the organization of labor and assets under an ownership nucleus that
retains residual control rights and captures residual output. In
Coase's analysis, the subordination layer is not incidental; it is
load-bearing, because enforcement between non-subordinated
counterparties was costly in the pre-primitive regime. A fixed-cost
bilateral-bond settlement primitive prices that enforcement
directly. The subordination layer loses its cost advantage for the
class of transactions identified below; what emerges is not the
abolition of the firm, but its demotion from the privileged
organizational container to one pattern among several. The
political-economy diagnostic develops what this demotion looks
like from the outside; the present paper is internal
institutional-economics.
\end{abstract}

\medskip
\noindent\textbf{Keywords:}
transaction cost economics, theory of the firm, subordination,
coordination economics, organizational substrate, institutional
economics

% ════════════════════════════════════════════════════════════════════
\section{Introduction}
\label{sec:intro}

The boundary of the firm has been the central object of theoretical
attention in organizational economics since
\citet{coase1937}. Williamson's Transaction Cost
Economics~\citep{williamson1985} identified asset specificity,
uncertainty, and frequency as the conditions under which
hierarchical governance dominates market exchange. The
property-rights theory~\citep{hart1995} located the firm boundary at
the allocation of residual control rights over non-contractible
decisions. Across these traditions, a common premise is that market
contracting between strangers is expensive: search is costly,
negotiation is costly, and enforcement---through courts,
arbitrators, or platforms---is costly, slow, and
jurisdiction-dependent.

This paper asks what happens to that boundary when one of the
underlying costs---contract enforcement between mutually distrusting
parties---collapses to a fixed, jurisdiction-displaced capital
lockup. The collapse is not a thought experiment: the mechanism
admits formal proofs of cooperation as the unique surviving Nash
equilibrium, and the deployed Solidity implementation has been
audited through TLA\textsuperscript{+} model checking,
property-based fuzzing,
symbolic execution, and specification checking.

We treat that primitive as exogenous. The contribution of this paper
is institutional-economic: given the primitive, what follows for the
theory of the firm? We argue two things. First, the Coasean
threshold shifts inward --- narrowly, in a precise sense we make
explicit. Second, the firm ceases to be the privileged organizational
container in the region where the threshold has shifted; it becomes
one pattern among several, sharing the organizational space with
treasury communities, peer networks, transaction-scoped assemblies,
agent coalitions, and forms whose naming is yet to stabilize.

\paragraph{Paper organization.}
\Cref{sec:primitive} summarizes the settlement primitive,
taking the formal development and implementation as given.
\Cref{sec:threshold} develops the Coasean threshold shift.
\Cref{sec:patterns} treats the firm as one organizational pattern
among many, introducing subordination as the variable whose cost
advantage the primitive erases, and developing the three concrete
implications for property-rights theory and for the unit of
coordination.
\Cref{sec:token} examines token denomination as a coordination
signal that supplants standing platform credentialing.
\Cref{sec:conclusion} concludes.

% ════════════════════════════════════════════════════════════════════
\section{The Settlement Primitive}
\label{sec:primitive}

The settlement primitive composes two distinct mechanisms that
together do the institutional work this paper draws on.
\emph{Asymmetric bonding} produces the bilateral Nash equilibrium
and scales it to $N$-party process trees through progressive
collateralization: each party locks capital against the
transaction's cumulative value, and cooperation is the unique
profile surviving iterated elimination of weakly dominated
strategies at every chain position. \emph{Buyer dominance with
atomic resolution} operates on the already-scaled mesh: only the
buyer can trigger resolution, and resolution settles all orders in
a process simultaneously or not at all. The two mechanisms are
inseparable: bonding alone yields independently secured bilateral
edges that cannot multi-party coordinate without external
choreography; buyer dominance alone is without force absent the
bonding equilibrium. Together they make the mesh resolvable from a
single signature with cooperation pressure propagating through it.
We use ``Mechanism 1'' and ``Mechanism 2'' below where the
distinction is load-bearing.

We summarize here only the properties of the primitive that the
present argument uses. The full definitions, theorems,
implementation, and verification belong to their own registers
and are taken as given here.

\paragraph{Asymmetric bonding.}
For a transaction with payment $\pay > 0$ and cumulative upstream
value $\cumval \geq \pay$, the buyer locks $2\pay$ and the seller
locks $2\cumval$. Only the buyer can trigger resolution. On
resolution, the buyer recovers $\pay$ (net cost: $-\pay$) and the
seller recovers $2\cumval + \pay$ (net gain: $+\pay$). On
non-resolution, both bonds remain locked indefinitely.

\paragraph{Equilibrium properties.}
In the resulting one-shot bonding game, cooperation weakly dominates
defection for all parties, and the cooperative profile is the unique
outcome surviving iterated elimination of weakly dominated
strategies. The result extends to $N$-party process trees through
\emph{progressive collateralization}, with the seller's
risk-to-reward ratio growing with chain depth.

\paragraph{Atomic resolution and peer enforcement.}
All orders within a process resolve atomically---all or
nothing---which induces a weakest-link subgame among sellers.
This produces an endogenous coordination pressure analogous to
peer enforcement in joint-liability microfinance, but obtained
without repeated interaction, local information, or social
punishment technology.

\paragraph{No escape hatches.}
The primitive has no timeout, no admin override, no governance
mechanism, no protocol fee, and no upgrade path. The
mechanism-design analysis proves that any such escape hatch
weakens the equilibrium.

\paragraph{What the primitive does \emph{not} provide.}
For the institutional argument that follows, two non-properties
matter as much as the properties. The primitive does not provide
capital aggregation across projects, does not provide pooled
risk-bearing across uncorrelated ventures, and does not provide
legal personhood for off-chain liability. We return to these
absences in \Cref{sec:patterns}.

% ════════════════════════════════════════════════════════════════════
\section{The Coasean Threshold Shift}
\label{sec:threshold}

Coase's central insight is that firms exist because market
transaction costs---searching for partners, negotiating terms,
enforcing agreements---sometimes exceed the cost of hierarchical
management. Williamson refined this into TCE, identifying three
conditions under which hierarchical governance dominates market
exchange: high asset specificity, high uncertainty, and high
frequency.

The settlement primitive changes the calculus by collapsing one of
these costs to a fixed, predictable capital lockup. Consider the
three transaction costs separately.

\paragraph{Search costs.}
Public on-chain signals (settlement history, accepted-token lists,
attested capabilities) allow autonomous discovery through shared
coordination graphs. For agents able to read blockchain state,
search cost approaches zero. The substantive change is not that
search is cheaper---platforms already make search cheap---but that
the search index is non-proprietary: no platform owns the discovery
function, and the data is not extracted as rent.

\paragraph{Negotiation costs.}
Off-chain commitment building (EIP-712 typed structured data) lets
parties iterate on terms without gas costs. The negotiation is
bilateral and direct; counterparties can discover and iterate
through shared graph visibility rather than through a standing
intermediary's matching function. The substantive change here is
asymmetric: the \emph{cost} of negotiating is roughly unchanged from
web2 baselines, but the \emph{counterparty space} from which one
can negotiate is no longer mediated by a platform's terms of
service.

\paragraph{Enforcement costs.}
This is the decisive variable. In the pre-primitive world,
enforcement requires courts (slow, expensive,
jurisdiction-dependent), arbitrators (trusted third parties), or
standing intermediaries that internalize coordination and charge
for it. The primitive prices enforcement at exactly $2\times$ the
transaction value---a fixed, predictable, jurisdiction-displaced
cost requiring no third party.\footnote{The $2\pay$ figure is the
capital-lockup component, which is the substantive change relative
to the prior regime. A full enforcement-cost accounting adds
(i) gas for the \texttt{commit} and \texttt{resolveProcess} calls
(bounded but non-zero), (ii) opportunity cost on locked capital
over the process duration $T$ at the prevailing risk-free rate $r$
(roughly $2\pay \cdot rT$), (iii) slippage if bond-denomination
currency must be acquired at transaction time, and
(iv) operational-security costs associated with key custody
(signing infrastructure, backup, rotation). The Coasean-threshold
comparison below abstracts (i)--(iv) because they are small
relative to the standing-firm overheads against which the
comparison is made for the retail-scale transactions this paper
treats; the abstraction is not universally valid and should be
revisited for applications in which any of (i)--(iv) is
comparable in magnitude to $2\pay$.}

\begin{proposition}[Coasean Threshold Shift]
\label{prop:coasean-shift}
For any transaction where $2\pay$ (the capital lockup cost) is
less than the firm's coordination overhead for the same
transaction, the rational organizational choice is market exchange
via the primitive rather than hierarchical integration.
\end{proposition}

\Cref{prop:coasean-shift} is a comparative-statics statement, not a
sweeping institutional prediction. It identifies a class of
transactions for which the make-or-buy boundary tilts toward
\emph{buy}---specifically, transactions where the time-weighted
opportunity cost of locked capital is lower than the standing
overhead of internalizing the coordination function. The class is
non-empty (most retail-scale transactions in mature financial
environments satisfy it) and bounded (very-high-frequency,
very-low-margin, or very-long-duration transactions do not). The
proposition does not claim that all transactions migrate; it claims
that the migration threshold is now defined by capital opportunity
cost rather than by the cost of standing institutional infrastructure.

\paragraph{Engagement with TCE.}
Williamson's three drivers of vertical integration deserve direct
treatment. \emph{Asset specificity} continues to favor hierarchy: a
specialized asset whose value is contingent on a specific
counterparty cannot be efficiently coordinated through a sequence of
short-lived bonded transactions, since the holdup risk is in the
asset's pre-transaction sunk investment, not in the transaction
itself. \emph{Uncertainty} cuts both ways: contract incompleteness
favors hierarchy, but the primitive's evidence trail (timestamped
commitments, signed disclosures) reduces some classes of
uncertainty---specifically those that hierarchical authority
historically resolved by managerial fiat. \emph{Frequency} is the
condition under which the shift is most pronounced: high-frequency
exchanges with low asset specificity are exactly where the standing
firm's overhead is hardest to amortize and the primitive's
per-transaction capital lockup is easiest to absorb.

The proposition therefore narrows rather than overturns TCE. It
identifies a region of the asset-specificity / uncertainty /
frequency space where hierarchical governance was previously the
equilibrium outcome but is no longer. \citet{tadelis_williamson_2013}
update this picture for the contemporary literature, treating the
transaction as the unit of analysis and the choice among market,
hybrid, and hierarchy as turning on transaction attributes;
\Cref{prop:coasean-shift} adds a new entry to that menu of
governance modes --- the bonded transaction-scoped assembly --- whose
attributes are different from any of the three Williamson catalogs.

\paragraph{Holmstrom and Roberts' caution.}
\citet{holmstrom_roberts_1998} argue that asset-specificity-based
hold-up is not by itself sufficient to explain firm boundaries:
firms also do work that markets do not, particularly in promoting
cooperation and coordinating decisions that no single contract can
fully specify. The caution applies to our argument. Mechanism~1
prices the enforcement cost; that addresses one component of the
firm's bundle, not all of it. The cooperation-promotion work that
Holmstrom and Roberts attribute to firms is precisely what
Mechanism~2 supplies through the weakest-link subgame: each
seller's payout depends on universal cooperation, so the
cooperation-promotion function is
discharged by the bonding mesh's structure rather than by a
manager. The two-mechanism composition addresses both components
of the firm's bundle that the Holmstrom--Roberts critique
identifies. The non-coordination components --- capital aggregation,
risk pooling, organizational knowledge --- remain durable
requirements, and we return to them in
\Cref{sec:patterns}.

% ════════════════════════════════════════════════════════════════════
\section{The Firm as One Pattern Among Many}
\label{sec:patterns}

The implication of \Cref{prop:coasean-shift} is not that the firm
dissolves but that the firm's status in the organizational space
changes. The firm has always been a coordination mechanism. What
distinguishes the modern corporate firm from older coordination
forms---partnerships, guilds, cooperatives, market exchange---is not
the coordination function itself but a particular overlay placed on
top of it: \emph{subordination}.

\subsection{Subordination as the Variable}

By subordination we mean the allocation of labor and productive
assets under a single ownership nucleus, combined with an
employment contract that transfers residual output rights and
residual control rights to that nucleus in exchange for a wage. In
the Coasean analysis, subordination is not incidental to the firm;
it is load-bearing. Prior to the settlement primitive, enforcement
between non-subordinated counterparties was costly, and aggregation
under an ownership nucleus that resolved enforcement internally
through managerial authority was the institutional response.

\begin{proposition}[Subordination as an Enforcement Economy]
\label{prop:subordination}
In the pre-primitive regime, subordination is a response to
enforcement cost. The firm's subordination layer secures
coordination between counterparties whose direct bilateral
enforcement would be prohibitively costly, at the price of
aggregating residual rights under a single nucleus. When bilateral
enforcement is priced at a fixed capital lockup, the subordination
layer loses its cost advantage in the region identified by
\Cref{prop:coasean-shift}.
\end{proposition}

\Cref{prop:subordination} reframes what is dissolving. It is not
the firm's coordination function that becomes unnecessary ---
coordination is still required. It is the subordination overlay ---
the transfer of residual control rights from the producing party to
an ownership nucleus --- that becomes economically optional where
the primitive is applicable. What is left behind is coordination
without subordination: the same productive assets, the same
coordination requirements, but without the ownership aggregation
that historically secured them.

\subsection{What Emerges in the Shifted Region}

The viability of the alternative patterns named below rests on a
property of the primitive that \Cref{sec:primitive} introduced but
that the Coasean argument of \Cref{sec:threshold} did not yet
deploy: atomic resolution. Mechanism~1 (asymmetric bonding)
supplies the bilateral equilibrium at every edge of an arbitrary
process tree, but a mesh of independently secured bilateral edges
is not by itself an organization --- coordinating an
$N$-party process to a single settlement under Mechanism~1 alone
would still require either a standing intermediary, an arbitrator,
or an $N$-way mutual signing at the moment of resolution. Any of
those three responses re-introduces the firm-shaped container
under a different name. Mechanism~2 (buyer dominance with atomic
resolution) makes the mesh of edges resolvable as a unit at the
buyer's signature, and the weakest-link subgame propagates
cooperation pressure through the mesh without managerial
authority or external choreography.
The transaction-scoped assembly clears itself. The Coasean shift
of \Cref{sec:threshold} is what makes the assembly competitive
with the firm at the cost margin; Mechanism~2 is what makes the
assembly structurally available at all.

Where subordination is no longer load-bearing, the firm becomes one
organizational pattern among several. The organizational space in
that region admits multiple viable coordination patterns that were
previously suppressed by the enforcement-cost structure:
transaction-scoped assemblies that form around a process and
dissolve at settlement; standing treasury communities whose members
recurrently participate in compatible process trees; peer networks
of independent bonded counterparties without a common treasury;
agent coalitions mixing human, software, and machine participants.
The firm continues to exist; it shares the space with these other
forms rather than occupying it alone. Which form dominates for a
given industry is an empirical question about the parameters of
\Cref{prop:coasean-shift} and the non-coordination functions
(capital aggregation, risk pooling, legal personhood) that the
primitive does not itself address.

\paragraph{Hansmann's typology of ownership.}
The closest comparator in the existing literature is
\citet{hansmann1996ownership}, who classifies firms by the patron
class that holds residual claims and control rights: investor-owned,
worker-owned, customer-owned, supplier-owned, and mutual. Each
form is an ownership nucleus, and Hansmann's central observation is
that one form dominates within an industry when the costs of
contracting with the other patron classes exceed the costs of
ownership by the dominant class. Transaction-scoped assemblies sit
outside this typology: they have no standing residual claim to
allocate, because resolution dissolves the per-process residual
into deterministic per-order payouts (the conservation law of the
mechanism paper). This is not a fifth Hansmann form; it is what
becomes available when the residual-claim allocation question is
discharged at settlement rather than carried over time. Hansmann's
heterogeneity-cost objection to worker cooperatives --- that
collective decision-making among workers with different preferences
is expensive enough to outweigh the benefits of worker ownership ---
does not bind here because the assembly's lifespan is bounded by a
single process: there is no continuing collective decision to make.
Where the assembly form is feasible and the firm form is not (the
shifted region of \Cref{prop:coasean-shift}), the assembly is
chosen for the same Hansmann-style reason that a particular firm
form is chosen elsewhere: the contracting costs against the
alternative are higher.

\paragraph{Ostrom's design principles.}
\citet{ostrom1990governing} identifies eight design principles
that long-enduring self-governing institutions tend to share:
clearly defined boundaries, congruence between rules and local
conditions, collective-choice arrangements, monitoring, graduated
sanctions, conflict-resolution mechanisms, recognition by external
authorities, and (for systems above a threshold size) nested
enterprises. Transaction-scoped assemblies satisfy several of
these structurally rather than by governance: boundaries are
defined by who has signed a process commitment; sanctions are
graduated by the bonded mesh's payoff structure (each seller's
bond is forfeit, not penalty-as-a-policy); conflict resolution
defaults to atomic resolution at the buyer's signature, with
external legal forums available as evidentiary consumers of the
bonded commitment. Other principles --- collective
choice, congruence with local conditions --- are properties of the
graph the parties compose on top of the kernel, not of the kernel
itself. The point is not that the assembly is an Ostrom commons
(it is not); it is that several Ostrom-identified properties of
durable institutions appear here as architectural rather than as
governance achievements.

\paragraph{Relational contracts and \citeauthor{macneil1980new}.}
The relational-contract literature \citep{macneil1980new,
gibbons2005four} treats long-running counterparty relationships
as governance structures in their own right --- self-enforcing
through reputation and continuation value when formal contracting
is incomplete. The bonded primitive does not replace relational
contracting; it composes with it. A standing peer network or
treasury community is precisely a long-running relationship in
\citeauthor{macneil1980new}'s sense, and the primitive supplies
the per-transaction enforcement that the relational contract
would otherwise have to underwrite via reputation. The two
substitute for each other along one axis (formal enforcement of
each transaction) and complement each other along another (the
relational layer carries reputation; the bonded layer carries
the equilibrium).

\paragraph{Holacracy as an intra-firm precedent.}
Intra-firm experiments with flat coordination --- most visibly the
Zappos adoption of Holacracy~\citep{robertson2015holacracy} in 2014
--- previewed what some of these patterns look like when they
remain confined to a single ownership nucleus. Such experiments are informative as data points
about what flat coordination requires operationally, but they are
not what the primitive generalizes. Holacracy, once the
subordination nucleus is retained, operates at the mercy of that
nucleus. The primitive enables flat coordination where the nucleus
itself is not required; that is a stronger condition than any
intra-firm reform can provide.

\subsection{Three Concrete Implications}

\paragraph{(a) The wallet is the unit of coordination.}
Each productive asset is represented directly by a wallet address.
At the protocol tier\footnote{We use ``kernel / protocol /
runtime'' in the sense defined in §1.4 of the mechanism paper.
The kernel guarantees only the \texttt{commit}/\texttt{resolveProcess}
bond structure; the credential and attestation apparatus described
in this paragraph is realized at the protocol tier by extensions
(attestation coordinators, schema registries) whose composition
discipline is part of the implementation surface.} the wallet can hold
content-addressed credentials --- attested capabilities, operator
profiles, and (at the runtime tier) NFT-anchored deeds --- that
certify the asset's identity and performance history. In the firm
model, the cook's skill, the kitchen's equipment, the courier's
vehicle are subsumed under the firm's legal identity and cease to be
individually addressable. In the wallet model, each asset is its
own coordination endpoint. The productive asset \emph{is} the
identity; it remains an identity whether the firm persists or not.

\paragraph{(b) Tokens replace stock.}
Participants in the shifted region do not receive shares in a firm;
they hold tokens denominated by the communities with which they
coordinate. Token holdings signal community affiliation and
participation-history, not residual-output ownership in a
centralized nucleus. The relationship between productive asset and
community is market-mediated rather than hierarchically
internalized. Tokens are borderless, community-defined rather than
jurisdiction-issued, and the community's definition is itself an
open parameter.

\paragraph{(c) Actor type is neutral.}
Nothing in the primitive's coordination layer distinguishes an AI
capable of bearing capital risk from a human capable of the same
risk. A person with a law degree and an AI trained on legal
precedent compete on identical terms --- each as a bonded
counterparty whose capabilities, if attested at all, are attested
at the protocol tier by a schema-validated attestation coordinator
rather than asserted at the kernel tier. Actor-type neutrality is
a property of the kernel's coordination layer, not a policy
choice; it follows directly from the kernel's indifference to what
sits behind a wallet.

\subsection{Engagement with the Property-Rights Tradition}

The property-rights tradition~\citep{hart1995, grossman_hart_1986,
hart_moore_1990} locates the firm boundary at the allocation of
residual control rights over non-contractible decisions. The
primitive does not abolish non-contractibility; rather, it changes
\emph{where} residual control rights are allocated, and it
narrows the temporal window over which they are exercised. In the
firm model, those rights are allocated to the ownership nucleus and
exercised continuously through managerial authority. In the wallet
model, they are allocated to whoever controls the wallet ---
which, under self-sovereign key management and NFT-anchored
credentials, is the productive asset itself (or the individual
operating it) --- and Mechanism~2 closes the temporal window:
atomic resolution dissolves the per-process residual claim into
deterministic per-order payouts at settlement, leaving no
ongoing residual that requires standing custody. Residual control
rights do not disappear; they are held by the asset rather than
aggregated under an ownership layer, and the per-process residual
claim is discharged at the moment of resolution rather than
accumulated over time.

This is the institutional-economics translation of what
\Cref{sec:patterns} has been describing throughout: the decision
about who holds residual rights, previously answered by the firm's
ownership structure, is now a property of the asset's wallet
architecture. This relocates the allocation question rather than
dissolving it; \citet{hart_moore_1990} applies to the wallet-layer
question just as it applied to the firm-layer question before,
with a different set of empirical facts about costs and
observability.

\paragraph{Williamson's drivers, revisited.}
Our treatment of Williamson's three drivers
(\Cref{sec:threshold}) likewise does not dissolve TCE; it relocates
the question. Asset specificity continues to matter wherever a
productive asset's value is contingent on a specific bond-specific
investment; the primitive does not change that. What changes is
that the institutional response to asset specificity need not be
subordination under an ownership nucleus --- a treasury community,
a long-term peer network, or a serial bonded-commitment pattern
can, for some parameter regions, secure the specific investment
without aggregating residual rights.

The non-coordination functions of the firm --- capital aggregation,
organizational knowledge, risk pooling, legal personhood --- persist
as durable requirements that some institution must satisfy. They do
not necessarily re-bundle into a firm-shaped container; they can be
satisfied by treasuries, DAOs, mutuals, cooperatives, or
jurisdiction-specific legal wrappers attached to wallet clusters.
The empirical question of which bundling pattern dominates in a
given industry is outside the scope of this paper.

\paragraph{The firm's self-knowledge function dissolves alongside.}
A function we have not separately named is that the firm
historically has been the unit at which \emph{books} are kept ---
the unit whose position is tracked over time through
double-entry bookkeeping in the tradition of
\citet{pacioli1494summa}. The bookkeeping function is not
incidental: at scale, an entity that cannot maintain accurate
records of its position cannot grow, attract investment, pay
taxes, or contract with third parties who require disclosure.
A \figaro{} process is a self-closing ledger period --- and the
apparatus that grew up
around the books (bookkeeping, accounting, audit) becomes
optional for activity that settles through the protocol. The
implication for the present argument is that the firm's
demotion from privileged organizational container is supported
by an additional structural shift not addressed elsewhere
above: the records function that motivated the firm boundary
historically is also no longer firm-bound. A wallet cohort
participating in many process trees has its full position
visible without consolidation.

% ════════════════════════════════════════════════════════════════════
\section{Token Denomination as Coordination Signal}
\label{sec:token}

In the primitive, the choice of denomination token is itself a
coordination signal. The settlement layer is token-agnostic: any
ERC-20 can serve as the bond and payment currency. This means:

\begin{itemize}
  \item A stablecoin signals legal-system alignment.
  \item A DAO governance token signals community membership.
  \item A commodity-backed token signals value anchoring.
  \item Accepting a specific token \emph{is} the seller's identity
    declaration---an accepted-token list is a brand.
\end{itemize}

Settlement velocity in a given token replaces traditional
reputation metrics (star ratings, review scores). The on-chain
record of resolved orders denominated in token $\tau$ is a public,
tamper-proof signal of the seller's reliability within the $\tau$
economy. No platform needs to aggregate, curate, or attest this
signal; it is a natural byproduct of protocol usage.

This matters institutionally because reputation infrastructure is
one of the durable rents that platforms extract. A standing
platform's credentialing function---certifying that a counterparty
is reliable, safe, or qualified---is one of the structural reasons
platforms persist even after their matching function is technically
replicable. When reputation is a public byproduct of settlement
rather than a proprietary platform asset, the platform's claim to
that rent weakens.

\paragraph{Schelling-point tokens as a distinct category.}
The coordination-signal argument above applies to \emph{any}
token used in a bonded commitment. A second coordination role is
played by tokens designed specifically to serve as focal points
in the sense of \citet{schelling1960strategy}: anchors around
which participants converge to signal alignment with a specific
ecosystem, without central authority and without relying on any
of the token's financial-mechanism features. The FIG token
treatment develops this category in detail, with \textsc{Btc}
and \textsc{Eth} as precedents at the
money and settlement layers respectively, and the \fig{} token
as the focal anchor for participants who coordinate around the
bonded-commitment primitive specifically. A Schelling-point
token does not compete with the general-purpose denomination
tokens discussed above; it carries a signal they cannot carry,
which is alignment with a specific coordination ecosystem rather
than with legal-system, community, or value-anchoring choices.
Such a token is underpinned not by cost-of-production (as Bitcoin
is) nor by transaction-execution demand (as Ethereum is) but by
the \emph{value-added of trade} conducted through the substrate
it coordinates around: \fig{}'s value rests on the scale and
quality of bonded commitments conducted by participants who use
it to identify each other.

% ════════════════════════════════════════════════════════════════════
\section{Conclusion}
\label{sec:conclusion}

The Coasean threshold is not a conjecture---it is observable.
Contemporary coordination stacks often extract 15--30\% of
transaction value as coordination rent (Uber~25\%, DoorDash~30\%,
Airbnb~14\%, Amazon Marketplace~15\%). Credit card networks layer
2--3\% interchange plus double-digit annual interest on revolving
balances. Financial markets intermediate trillions in notional
value against a fraction of that in real economic exchange.
Asymmetric bonding (Mechanism~1) prices the cost of trust at
$2\times$ the transaction value in temporarily locked (not spent)
capital, with zero recurring extraction. Buyer dominance with
atomic resolution (Mechanism~2) makes the resulting mesh of
bonded edges resolvable as a single coordinated unit at the buyer's
signature. The Coasean threshold shift is the work of Mechanism~1;
the structural availability of transaction-scoped assemblies as a
firm-substituting organizational form is the work of Mechanism~2.

For any economic activity where the coordination rent exceeds the
opportunity cost of locked capital, the standing firm --- as Coase
conceived it --- is no longer the uniquely efficient unit of
economic organization in that activity. The stronger consequence is
that the firm becomes one organizational pattern among several in
the shifted region: transaction-scoped assemblies, treasury
communities, peer networks, and agent coalitions share the
coordination space alongside it. The settlement primitive does not
abolish coordination, nor does it abolish the firm; it retires the
\emph{subordination overlay} that previously made the firm the
privileged container, and leaves a richer organizational space in
its place.

The argument leaves three questions open. First, the empirical
question: which industries and which transaction classes lie above
the threshold today, and at what rate does the threshold shift as
on-chain settlement matures? Second, the institutional question:
which organizational patterns --- firms, treasuries, peer networks,
or forms not yet named --- dominate in a given industry once the
primitive is available, and what determines the dominance? The
non-coordination functions of the firm (capital aggregation,
organizational knowledge, risk pooling, legal personhood) must
still be bundled somewhere; the bundling question is open. Third, a
question of a different character: why this shift is qualitatively
harder to recognize than to specify, and what that difficulty
itself reveals about the architecture of contemporary
coordination infrastructure. That third question belongs to the
political-economy register and is treated there.

\paragraph{A note on scope.}
This paper bounds its claim by two scope conditions: the Coasean
comparison (\Cref{prop:coasean-shift}: $2\pay$ vs. coordination
overhead) and asset specificity (\Cref{sec:patterns}). Both must
hold for a transaction to lie within the paper's argumentative
domain; readers constructing arguments about transactions that
satisfy one condition but not the other should treat the
argument with corresponding partial force. Adjacent diagnostic
treatments bound their claims by different conditions (hegemonic
legibility, labor-law specificity, evidentiary access); the
present paper does not adjudicate among those bounds.

\section*{Acknowledgments}
I thank the readers and collaborators who pressed me to separate
the institutional-economics argument made here from its
political-economy half (now treated separately). The two were
entangled in earlier drafts in a way that did neither audience
justice.


% ════════════════════════════════════════════════════════════════════
% References
% ════════════════════════════════════════════════════════════════════

\begin{thebibliography}{99}

\bibitem[Coase(1937)]{coase1937}
Coase, R.~H.
\newblock The Nature of the Firm.
\newblock \emph{Economica}, 4(16):386--405, 1937.

\bibitem[Pacioli(1494)]{pacioli1494summa}
Pacioli, L.
\newblock \emph{Summa de Arithmetica, Geometria, Proportioni
  et Proportionalita}.
\newblock Paganinus de Paganinis, Venice, 1494.

\bibitem[Grossman and Hart(1986)]{grossman_hart_1986}
Grossman, S.~J. and Hart, O.~D.
\newblock The Costs and Benefits of Ownership: A Theory of Vertical
  and Lateral Integration.
\newblock \emph{Journal of Political Economy}, 94(4):691--719, 1986.

\bibitem[Hart(1995)]{hart1995}
Hart, O.
\newblock \emph{Firms, Contracts, and Financial Structure}.
\newblock Oxford University Press, 1995.

\bibitem[Hart and Moore(1990)]{hart_moore_1990}
Hart, O. and Moore, J.
\newblock Property Rights and the Nature of the Firm.
\newblock \emph{Journal of Political Economy}, 98(6):1119--1158, 1990.

\bibitem[Robertson(2015)]{robertson2015holacracy}
Robertson, B.~J.
\newblock \emph{Holacracy: The New Management System for a Rapidly
  Changing World}.
\newblock Henry Holt and Company, New York, 2015.

\bibitem[Schelling(1960)]{schelling1960strategy}
Schelling, T.~C.
\newblock \emph{The Strategy of Conflict}.
\newblock Harvard University Press, Cambridge, MA, 1960.

\bibitem[Williamson(1985)]{williamson1985}
Williamson, O.~E.
\newblock \emph{The Economic Institutions of Capitalism}.
\newblock Free Press, New York, 1985.

\bibitem[Holmstrom and Roberts(1998)]{holmstrom_roberts_1998}
Holmstrom, B. and Roberts, J.
\newblock The Boundaries of the Firm Revisited.
\newblock \emph{Journal of Economic Perspectives}, 12(4):73--94, 1998.

\bibitem[Hansmann(1996)]{hansmann1996ownership}
Hansmann, H.
\newblock \emph{The Ownership of Enterprise}.
\newblock Harvard University Press, Cambridge, MA, 1996.

\bibitem[Ostrom(1990)]{ostrom1990governing}
Ostrom, E.
\newblock \emph{Governing the Commons: The Evolution of Institutions
  for Collective Action}.
\newblock Cambridge University Press, Cambridge, 1990.

\bibitem[Gibbons(2005)]{gibbons2005four}
Gibbons, R.
\newblock Four Formal(izable) Theories of the Firm?
\newblock \emph{Journal of Economic Behavior \& Organization},
  58(2):200--245, 2005.

\bibitem[Tadelis and Williamson(2013)]{tadelis_williamson_2013}
Tadelis, S. and Williamson, O.~E.
\newblock Transaction Cost Economics.
\newblock In R.~Gibbons and J.~Roberts, editors, \emph{Handbook of
  Organizational Economics}, chapter~4. Princeton University Press,
  Princeton, NJ, 2013.

\bibitem[Macneil(1980)]{macneil1980new}
Macneil, I.~R.
\newblock \emph{The New Social Contract: An Inquiry into Modern
  Contractual Relations}.
\newblock Yale University Press, New Haven, 1980.

\end{thebibliography}

\end{document}
